\documentclass[10pt]{extarticle}
\usepackage[a4paper,margin=1.5cm]{geometry}
\usepackage[T1]{fontenc}
\usepackage{amsmath,amssymb}
\usepackage{graphicx}
\usepackage[export]{adjustbox}
\usepackage{multicol}
\usepackage{enumitem}
\usepackage{booktabs,array,tabularx}
\setlength{\columnsep}{0.7cm}
\setlength{\parindent}{0pt}
\setlength{\parskip}{3pt}
\pagestyle{empty}
\makeatletter
\newcommand{\qfigdef}[4]{\@namedef{qf@file@#1}{#2}\@namedef{qf@w@#1}{#3}\@namedef{qf@wide@#1}{#4}}
\newcommand{\qf@check}[1]{\@ifundefined{qf@file@#1}{\PackageError{score2}{Unknown figure id #1}{}}{}}
\newcommand{\qf@img}[4][]{\includegraphics[width=#2\dimexpr\@nameuse{qf@w@#4}pt\relax,max width=#3,max height=0.45\textheight,#1]{\@nameuse{qf@file@#4}}}
\newif\ifqf@full
\newif\ifqf@boxed
\newsavebox{\qf@box}
\newcommand{\qf@open}{\global\qf@fulltrue\par\global\qf@boxedtrue
  \setbox\qf@box\vbox\bgroup\hsize\textwidth\linewidth\textwidth}
\newcommand{\qf@emit}[1]{\par\egroup\global\qf@boxedfalse
  \ifdim\dimexpr\ht\qf@box+\dp\qf@box\relax<#1\textheight\noindent\usebox{\qf@box}\par\else\unvbox\qf@box\fi}
\newcommand{\qfullsplit}{\ifqf@boxed\qf@emit{0.85}\fi}
\newcommand{\qf@close}{\ifqf@boxed\qf@emit{0.45}\fi\global\qf@fullfalse}
\newcommand{\qfullbegin}{\end{multicols}\qf@open}
\newcommand{\qfullend}{\qf@close\begin{multicols}{2}}
\newcommand{\qfullstart}{\qf@open}
\newcommand{\qfulldone}{\qf@close}
\newcommand{\qfig}[2][1]{\qf@check{#2}%
  \ifqf@full\qf@blk{#1}{#2}\else
  \ifnum\@nameuse{qf@wide@#2}=1 \expandafter\@firstoftwo\else\expandafter\@secondoftwo\fi
  {\end{multicols}\begin{center}\qf@img{#1}{\textwidth}{#2}\end{center}\begin{multicols}{2}}%
  {\qf@blk{#1}{#2}}\fi}
\newcommand{\qf@blk}[2]{\par\smallskip{\centering\qf@img{#1}{\linewidth}{#2}\par}\smallskip}
\newcommand{\qfigi}[2][1]{\qf@check{#2}\qf@img[valign=c]{#1}{\linewidth}{#2}}
\makeatother
\newcommand{\papersec}[1]{\par\bigskip{\centering\Large\bfseries PAPER-#1\par}\medskip}
\newcommand{\partsec}[2]{\par\medskip\noindent{\large\bfseries PART-#1: #2}\par\smallskip}
\newcommand{\sect}[1]{\par\smallskip\noindent\textbf{#1}\par}
\newcommand{\qn}[1]{\par\smallskip\noindent\textbf{#1.}~}
\newcommand{\paraq}[2]{\par\smallskip\noindent\textbf{Paragraph for Questions #1 and #2}\par}
\newenvironment{qoptions}{\begin{itemize}[leftmargin=1.8em,itemsep=1pt,topsep=2pt,label={}]}{\end{itemize}}
\newcommand{\opt}[2]{\item[(#1)] #2}
% figure registry (generated by score2_tools.py finalize)
\qfigdef{p1_phy_s1_q01_1}{JEE Advanced (2021 Series) - Test 08 (01 May 2021)-img/p1_phy_s1_q01_1.png}{108.5}{0}
\qfigdef{p1_phy_s1_q05_1}{JEE Advanced (2021 Series) - Test 08 (01 May 2021)-img/p1_phy_s1_q05_1.png}{108.5}{0}
\qfigdef{p1_phy_s1_q07_1}{JEE Advanced (2021 Series) - Test 08 (01 May 2021)-img/p1_phy_s1_q07_1.png}{78.7}{0}
\qfigdef{p1_phy_s2_para11_13_1}{JEE Advanced (2021 Series) - Test 08 (01 May 2021)-img/p1_phy_s2_para11_13_1.png}{199.1}{0}
\qfigdef{p1_phy_s3_q01_1}{JEE Advanced (2021 Series) - Test 08 (01 May 2021)-img/p1_phy_s3_q01_1.png}{100.0}{0}
\qfigdef{p1_phy_s3_q03_1}{JEE Advanced (2021 Series) - Test 08 (01 May 2021)-img/p1_phy_s3_q03_1.png}{67.3}{0}
\qfigdef{p1_chem_s1_q02_B}{JEE Advanced (2021 Series) - Test 08 (01 May 2021)-img/p1_chem_s1_q02_B.png}{52.0}{0}
\qfigdef{p1_chem_s1_q02_C}{JEE Advanced (2021 Series) - Test 08 (01 May 2021)-img/p1_chem_s1_q02_C.png}{45.5}{0}
\qfigdef{p1_chem_s1_q02_D}{JEE Advanced (2021 Series) - Test 08 (01 May 2021)-img/p1_chem_s1_q02_D.png}{51.4}{0}
\qfigdef{p1_chem_s1_q05_A}{JEE Advanced (2021 Series) - Test 08 (01 May 2021)-img/p1_chem_s1_q05_A.png}{99.2}{0}
\qfigdef{p1_chem_s1_q05_B}{JEE Advanced (2021 Series) - Test 08 (01 May 2021)-img/p1_chem_s1_q05_B.png}{104.8}{0}
\qfigdef{p1_chem_s1_q05_C}{JEE Advanced (2021 Series) - Test 08 (01 May 2021)-img/p1_chem_s1_q05_C.png}{180.4}{0}
\qfigdef{p1_chem_s1_q05_D}{JEE Advanced (2021 Series) - Test 08 (01 May 2021)-img/p1_chem_s1_q05_D.png}{197.0}{0}
\qfigdef{p1_chem_s1_q07_A}{JEE Advanced (2021 Series) - Test 08 (01 May 2021)-img/p1_chem_s1_q07_A.png}{98.2}{0}
\qfigdef{p1_chem_s1_q07_B}{JEE Advanced (2021 Series) - Test 08 (01 May 2021)-img/p1_chem_s1_q07_B.png}{91.3}{0}
\qfigdef{p1_chem_s3_q01_1}{JEE Advanced (2021 Series) - Test 08 (01 May 2021)-img/p1_chem_s3_q01_1.png}{121.8}{0}
\qfigdef{p1_chem_s3_q02_1}{JEE Advanced (2021 Series) - Test 08 (01 May 2021)-img/p1_chem_s3_q02_1.png}{171.1}{0}
\qfigdef{p1_chem_s3_q05_1}{JEE Advanced (2021 Series) - Test 08 (01 May 2021)-img/p1_chem_s3_q05_1.png}{245.2}{0}
\qfigdef{p2_phy_s1_q01_1}{JEE Advanced (2021 Series) - Test 08 (01 May 2021)-img/p2_phy_s1_q01_1.png}{99.7}{0}
\qfigdef{p2_phy_s1_q04_1}{JEE Advanced (2021 Series) - Test 08 (01 May 2021)-img/p2_phy_s1_q04_1.png}{63.5}{0}
\qfigdef{p2_phy_s1_q07_1}{JEE Advanced (2021 Series) - Test 08 (01 May 2021)-img/p2_phy_s1_q07_1.png}{110.1}{0}
\qfigdef{p2_phy_s2_q01_1}{JEE Advanced (2021 Series) - Test 08 (01 May 2021)-img/p2_phy_s2_q01_1.png}{86.2}{0}
\qfigdef{p2_phy_s2_q02_1}{JEE Advanced (2021 Series) - Test 08 (01 May 2021)-img/p2_phy_s2_q02_1.png}{91.3}{0}
\qfigdef{p2_phy_s2_q03_1}{JEE Advanced (2021 Series) - Test 08 (01 May 2021)-img/p2_phy_s2_q03_1.png}{152.3}{0}
\qfigdef{p2_phy_s2_q04_1}{JEE Advanced (2021 Series) - Test 08 (01 May 2021)-img/p2_phy_s2_q04_1.png}{104.9}{0}
\qfigdef{p2_phy_s2_q06_1}{JEE Advanced (2021 Series) - Test 08 (01 May 2021)-img/p2_phy_s2_q06_1.png}{97.5}{0}
\qfigdef{p2_phy_s3_para15_16_1}{JEE Advanced (2021 Series) - Test 08 (01 May 2021)-img/p2_phy_s3_para15_16_1.png}{130.0}{0}
\qfigdef{p2_phy_s3_para17_18_1}{JEE Advanced (2021 Series) - Test 08 (01 May 2021)-img/p2_phy_s3_para17_18_1.png}{163.1}{0}
\qfigdef{p2_chem_s1_q01_1}{JEE Advanced (2021 Series) - Test 08 (01 May 2021)-img/p2_chem_s1_q01_1.png}{236.6}{0}
\qfigdef{p2_chem_s1_q01_A}{JEE Advanced (2021 Series) - Test 08 (01 May 2021)-img/p2_chem_s1_q01_A.png}{39.4}{0}
\qfigdef{p2_chem_s1_q01_B}{JEE Advanced (2021 Series) - Test 08 (01 May 2021)-img/p2_chem_s1_q01_B.png}{42.7}{0}
\qfigdef{p2_chem_s1_q01_C}{JEE Advanced (2021 Series) - Test 08 (01 May 2021)-img/p2_chem_s1_q01_C.png}{41.8}{0}
\qfigdef{p2_chem_s1_q01_D}{JEE Advanced (2021 Series) - Test 08 (01 May 2021)-img/p2_chem_s1_q01_D.png}{42.3}{0}
\qfigdef{p2_chem_s1_q03_1}{JEE Advanced (2021 Series) - Test 08 (01 May 2021)-img/p2_chem_s1_q03_1.png}{262.7}{1}
\qfigdef{p2_chem_s1_q03_A}{JEE Advanced (2021 Series) - Test 08 (01 May 2021)-img/p2_chem_s1_q03_A.png}{65.1}{0}
\qfigdef{p2_chem_s1_q03_B}{JEE Advanced (2021 Series) - Test 08 (01 May 2021)-img/p2_chem_s1_q03_B.png}{64.2}{0}
\qfigdef{p2_chem_s1_q03_C}{JEE Advanced (2021 Series) - Test 08 (01 May 2021)-img/p2_chem_s1_q03_C.png}{64.1}{0}
\qfigdef{p2_chem_s1_q03_D}{JEE Advanced (2021 Series) - Test 08 (01 May 2021)-img/p2_chem_s1_q03_D.png}{64.6}{0}
\qfigdef{p2_chem_s1_q05_1}{JEE Advanced (2021 Series) - Test 08 (01 May 2021)-img/p2_chem_s1_q05_1.png}{146.9}{0}
\qfigdef{p2_chem_s2_q01_A}{JEE Advanced (2021 Series) - Test 08 (01 May 2021)-img/p2_chem_s2_q01_A.png}{63.9}{0}
\qfigdef{p2_chem_s2_q01_B}{JEE Advanced (2021 Series) - Test 08 (01 May 2021)-img/p2_chem_s2_q01_B.png}{106.3}{0}
\qfigdef{p2_chem_s2_q01_C}{JEE Advanced (2021 Series) - Test 08 (01 May 2021)-img/p2_chem_s2_q01_C.png}{63.7}{0}
\qfigdef{p2_chem_s2_q01_D}{JEE Advanced (2021 Series) - Test 08 (01 May 2021)-img/p2_chem_s2_q01_D.png}{71.1}{0}
\qfigdef{p2_chem_s2_q05_1}{JEE Advanced (2021 Series) - Test 08 (01 May 2021)-img/p2_chem_s2_q05_1.png}{137.4}{0}
\qfigdef{p2_chem_s2_q07_A}{JEE Advanced (2021 Series) - Test 08 (01 May 2021)-img/p2_chem_s2_q07_A.png}{84.6}{0}
\qfigdef{p2_chem_s2_q07_B}{JEE Advanced (2021 Series) - Test 08 (01 May 2021)-img/p2_chem_s2_q07_B.png}{84.6}{0}
\qfigdef{p2_chem_s2_q07_C}{JEE Advanced (2021 Series) - Test 08 (01 May 2021)-img/p2_chem_s2_q07_C.png}{84.6}{0}
\qfigdef{p2_chem_s3_para17_18_1}{JEE Advanced (2021 Series) - Test 08 (01 May 2021)-img/p2_chem_s3_para17_18_1.png}{245.2}{0}
\qfigdef{p2_chem_s3_q03_A}{JEE Advanced (2021 Series) - Test 08 (01 May 2021)-img/p2_chem_s3_q03_A.png}{20.3}{0}
\qfigdef{p2_chem_s3_q03_B}{JEE Advanced (2021 Series) - Test 08 (01 May 2021)-img/p2_chem_s3_q03_B.png}{78.5}{0}
\qfigdef{p2_chem_s3_q03_C}{JEE Advanced (2021 Series) - Test 08 (01 May 2021)-img/p2_chem_s3_q03_C.png}{43.6}{0}
\qfigdef{p2_chem_s3_q03_D}{JEE Advanced (2021 Series) - Test 08 (01 May 2021)-img/p2_chem_s3_q03_D.png}{37.7}{0}
\qfigdef{p2_chem_s3_q04_A}{JEE Advanced (2021 Series) - Test 08 (01 May 2021)-img/p2_chem_s3_q04_A.png}{66.8}{0}
\qfigdef{p2_chem_s3_q04_B}{JEE Advanced (2021 Series) - Test 08 (01 May 2021)-img/p2_chem_s3_q04_B.png}{55.4}{0}
\qfigdef{p2_chem_s3_q04_C}{JEE Advanced (2021 Series) - Test 08 (01 May 2021)-img/p2_chem_s3_q04_C.png}{78.6}{0}
\qfigdef{p2_chem_s3_q04_D}{JEE Advanced (2021 Series) - Test 08 (01 May 2021)-img/p2_chem_s3_q04_D.png}{78.5}{0}

\begin{document}
\begin{center}{\Large\bfseries JEE Advanced (2021 Series) --- Test 08 (01 May 2021)}\end{center}
\begin{multicols}{2}
\papersec{1}
\gdef\paraq#1#2{\par\smallskip\noindent\textbf{Paragraph for Questions #1 -- #2}\par}
\partsec{1}{PHYSICS}
\sect{SEC-1}
\qn{1} An electric dipole is fixed along $y$-axis with its centre at origin. At a given instant a charge particle is moving along $x$-axis then
\qfig{p1_phy_s1_q01_1}
\begin{qoptions}
\opt{A}{Power of electric force on charge particle is zero.}
\opt{B}{Acceleration and velocity are perpendicular}
\opt{C}{Time rate of change in velocity is zero.}
\opt{D}{If charge is free to move, it always moves along $x$-axis.}
\end{qoptions}
\qn{2} In a spherical capacitor, we have two concentric spherical shells, the inner one carrying a charge $Q$ and outer one carrying a charge of $-Q$. If the inner shell is displaced from the center without touching the outer shell,
\begin{qoptions}
\opt{A}{the capacitance of the capacitor will increase}
\opt{B}{the capacitance of the capacitor will remain same.}
\opt{C}{the energy of capacitor will decrease.}
\opt{D}{the potential difference between inner and outershell will increase.}
\end{qoptions}
\qn{3} Two toy cars with different masses originally at rest are pushed apart by an ideal spring and released. Which of the following statement(s) are true?
\begin{qoptions}
\opt{A}{Both toy cars will acquire equal but opposite momenta.}
\opt{B}{Both toy cars will acquire equal kinetic energy.}
\opt{C}{The more massive toy car will acquire less speed.}
\opt{D}{The smaller toy car will experience an acceleration of greater magnitude.}
\end{qoptions}
\qn{4} Which of the following statement(s) are true about a stationary wave ?
\begin{qoptions}
\opt{A}{No particle will move with same amplitude.}
\opt{B}{All particles between two successive nodes reach their extreme position together in phase.}
\opt{C}{Displacement and velocity nodes co-exist.}
\opt{D}{Velocity nodes and pressure nodes co-exist.}
\end{qoptions}
\qn{5} In order to reduce its diameter, a wire is pulled through a small hole in a metal plate. The wire is made of metal whose specific heat capacity is $400$ J kg$^{-1}\,^\circ$C$^{-1}$ and on emerging from the hole it has a mass per unit length of $5$ g per m. A steady force of $600$ N is applied on wire. If all the heat generated is retained in the wire, we can calculate the rise in temperature of the wire.
\qfig{p1_phy_s1_q05_1}
\begin{qoptions}
\opt{A}{Work done on one kilogram of wire is $120{,}000$ J}
\opt{B}{Assuming all the work done is converted into heat, the temperature rise of the wire is $300^\circ$C}
\opt{C}{$0.042$ kg wire is produced each second if it emerges from the hole at a speed of $8.4$ ms$^{-1}$.}
\opt{D}{$0.021$ kg wire is produced each second if it emerges from the hole at a speed of $8.4$ ms$^{-1}$.}
\end{qoptions}
\qn{6} A spherical solid body is dropped inside a vast expanse of viscous liquid of large depth and of coefficient of viscosity $\eta$. The density of the solid is greater than that of liquid. The time taken by the body to attain $90\%$ of the steady state velocity is dependent on :-
\begin{qoptions}
\opt{A}{Density of the liquid}
\opt{B}{Density of the solid}
\opt{C}{Diameter of the sphere}
\opt{D}{Coefficient of viscosity}
\end{qoptions}
\qn{7} A disc is given an initial angular velocity $\omega_0$ and is placed on rough horizontal surface as shown. The quantities which will not depend on the co-efficient of friction is/are :
\qfig{p1_phy_s1_q07_1}
\begin{qoptions}
\opt{A}{The time until pure rolling begins}
\opt{B}{The displacement of the disc until pure rolling begins.}
\opt{C}{The velocity when pure rolling begins.}
\opt{D}{The work done by the force of friction.}
\end{qoptions}
\sect{SEC-2}
\paraq{8}{10}
In a standard YDSE set up, a slab of thickness $t$ \& refractive index $\mu$ is placed in front of the upper slit. Other symbols have usual meaning.

\end{multicols}
\begin{center}\small\setlength{\tabcolsep}{4pt}\renewcommand{\arraystretch}{1.7}
\begin{tabular}{|c|p{3cm}|c|p{3cm}|c|p{3cm}|}\hline
 \multicolumn{2}{|c|}{Column--1} & \multicolumn{2}{c|}{Column--2} & \multicolumn{2}{c|}{Column--3}\\\hline
(I) & $D/d=10^{+3}$ & (i) & $t=1\,\mu$m, $\mu=1.6$ & (P) & $\lambda=400$ nm\\\hline
(II) & $D/d=2\times10^{+3}$ & (ii) & $t=1.5\,\mu$m, $\mu=1.3$ & (Q) & $\lambda=500$ nm\\\hline
(III) & $D/d=5\times10^{+4}$ & (iii) & $t=2\,\mu$m, $\mu=1.6$ & (R) & $\lambda=600$ nm\\\hline
(IV) & $D/d=4\times10^{+3}$ & (iv) & $t=2.5\,\mu$m, $\mu=1.4$ & (S) & $\lambda=700$ nm\\\hline
\end{tabular}
\end{center}
\begin{multicols}{2}
\qn{1} In which of the following case, the central maxima shifts by $1$ mm ?
\begin{qoptions}
\opt{A}{(II) (i) (P)}
\opt{B}{(III) (ii) (R)}
\opt{C}{(IV) (iii) (Q)}
\opt{D}{(I) (iv) (S)}
\end{qoptions}
\qn{2} In which of the following case, the central maxima shifts by $2$ fringe width ?
\begin{qoptions}
\opt{A}{(I) (i) (P)}
\opt{B}{(II) (ii) (Q)}
\opt{C}{(III) (iii) (R)}
\opt{D}{(IV) (iv) (S)}
\end{qoptions}
\qn{3} If the intensity at central maxima is $I_0$, in which of the following case intensity at the centre of the screen is zero ?
\begin{qoptions}
\opt{A}{(I) (i) (P)}
\opt{B}{(II) (ii) (Q)}
\opt{C}{(III) (iii) (R)}
\opt{D}{(IV) (iv) (S)}
\end{qoptions}
\paraq{11}{13}
Column-1 gives various circuit diagram with a black box which can contain either inductor or resistor. Column-2 lists time constant for the various circuit while column-3 lists the value of current ($i_1$) at either $t=0$ or $t\to\infty$. Initially before switching on inductor does not carry any energy.
\qfig{p1_phy_s2_para11_13_1}
\qn{4} If the black box consists of $2R$ (a resistor of resistance $2R$), then correct matching will be :
\begin{qoptions}
\opt{A}{(I) (iv) (Q)}
\opt{B}{(II) (i) (P)}
\opt{C}{(I) (iv) (R)}
\opt{D}{(IV) (iii) (P)}
\end{qoptions}
\qn{5} If the black box consists of $R$ (a resistor of resistance $R$), then correct matching will be :
\begin{qoptions}
\opt{A}{(I) (ii) (P)}
\opt{B}{(IV) (i) (Q)}
\opt{C}{(II) (iv) (S)}
\opt{D}{(II) (iii) (R)}
\end{qoptions}
\qn{6} If the black box consists of an inductor $L/2$, then correct matching will be :-
\begin{qoptions}
\opt{A}{(III) (ii) (R)}
\opt{B}{(III) (i) (R)}
\opt{C}{(III) (iii) (S)}
\opt{D}{(III) (i) (P)}
\end{qoptions}
\sect{SEC-3}
\qn{1} Figure shows infinitely long bent wires in the plane of paper and each carrying a current $i$ with $O$ as origin and $x$ and $y$ axis as shown, co-ordinates of points $P$ \& $Q$ are $(2\text{m},0)$ and $(6\text{m},0)$ respectively. Magnetic induction at $P$ is given by $\dfrac{\alpha\mu_0 i}{32\pi}$. Fill the value of $\alpha$ in OMR sheet.
\qfig{p1_phy_s3_q01_1}
\qn{2} X-ray are incident normally on a crystal of lattice constant $0.6$ nm. The first order reflection on diffraction from the crystal occurs at an angle of $30^\circ$. The wavelength of the X-ray used is $\alpha$ nm. Fill $10\alpha$ in OMR sheet.
\qn{3} In the figure the voltmeter shows reading of $20$ V. The battery has a voltage of $U_0=100$ V no internal resistance. Find the ratio of the current going through right part of the potentiometer to the current going through the voltmeter, if the jockey divides potentiometer into ratio $1:2$, and bigger resistance of the potentiometer is located to the right of jockey.
\qfig{p1_phy_s3_q03_1}
\qn{4} The count rate from $100$ cm$^3$ of a radioactive liquid is $c$. Some of liquid is now discarded. The count rate of the remaining liquid is found to be $\dfrac{c}{10}$ after three half-lives. The volume of the remaining liquid after three half lives in cm$^3$ is $(10\,n)$ cm$^3$. Calculate $n$.
\qn{5} In surveying, sometimes a value can be found for an inaccessible length $\ell$ (such as the height of a tall tree) by measuring other three lengths $\ell_1,\ell_2,\ell_3$ in terms of which $\ell=\dfrac{\ell_1\ell_2}{\ell_3}$. Suppose we perform such an experiment and obtain the following results (in feet) $\ell_1=200\pm2$, $\ell_2=5.5\pm0.1$, $\ell_3=10.0\pm0.4$. If in this experiment value of $\ell$ is expressed as $\ell=110\pm x$, find value of $x$ (in feet)
\partsec{2}{CHEMISTRY}
\sect{SEC-1}
\qn{1} At $300$ K, a $0.5$ mole sample of air ($80\%$ by mol of N$_2$ \& $20\%$ by mol of O$_2$) is put in $10$ litre cylinder piston system in which $5$ litre water is already present, then -\\
$k_H$ for N$_2=10^5$ atm\\
$k_H$ for O$_2=4\times10^4$ atm
\begin{qoptions}
\opt{A}{$X_{\mathrm{N}_2}$ in H$_2$O will be $1.9704\times10^{-5}$}
\opt{B}{$X_{\mathrm{O}_2}$ in H$_2$O will be $1.2315\times10^{-5}$}
\opt{C}{On increasing temperature value of Henry's constant increases}
\opt{D}{In H$_2$O, N$_2$(g) is more soluble as compare to O$_2$(g) at same pressure}
\end{qoptions}
\qn{2} Identify correct match for total number of products ?
\begin{qoptions}
\opt{A}{Optically active ketones with molecular formula [C$_5$H$_{10}$O] $\xrightarrow{\text{LiAlH}_4}$ Total six alcohol products are formed}
\opt{B}{\qfigi{p1_chem_s1_q02_B} $\xrightarrow[\Delta]{\text{Conc. H}_2\text{SO}_4}$ Five alkenes are formed during reaction}
\opt{C}{\qfigi{p1_chem_s1_q02_C} $\xrightarrow[(\mathrm{S_N}1)]{\text{EtOH}}$ Two ether products are formed}
\opt{D}{\qfigi{p1_chem_s1_q02_D} $\xrightarrow[\text{CCl}_4]{\text{HBr}}$ Total eight mono bromo major products are formed}
\end{qoptions}
\qn{3} Which of the following processes is/are involved in the extraction of Cu from sulphide ores of copper containing iron sulphide.
\begin{qoptions}
\opt{A}{FeS is more readily converted to the oxide than Cu$_2$S}
\opt{B}{Liq. Matte is placed in converter for carbon reduction}
\opt{C}{Cu$_2$S is partially converted to Cu$_2$O by air blasting and then to metallic copper}
\opt{D}{Silica is added to combine with FeO}
\end{qoptions}
\qn{4}
\begin{gather*}
E^\circ_{\mathrm{A^+/A}}=0.5\,\text{V};\ E^\circ_{\mathrm{B^+/B}}=0.2\,\text{V}\\
E^\circ_{\mathrm{C^+/C}}=-0.3\,\text{V};\ E^\circ_{\mathrm{D^+/D}}=-0.1\,\text{V}
\end{gather*}
Based on above data select correct statement under standard condition -
\begin{qoptions}
\opt{A}{Best oxidising agent is A$^+$}
\opt{B}{Best reducing agent is D}
\opt{C}{C can displace D from it's solution}
\opt{D}{In galvanic cell obtained by electrodes A \& B current flow from B to A in external circuit}
\end{qoptions}
\qn{5} Identify reactions correctly matched with their major product ?
\begin{qoptions}
\opt{A}{\qfigi{p1_chem_s1_q05_A}}
\opt{B}{\qfigi{p1_chem_s1_q05_B}}
\opt{C}{\qfigi{p1_chem_s1_q05_C}}
\opt{D}{\qfigi{p1_chem_s1_q05_D}}
\end{qoptions}
\qn{6} Select correct statement(s).
\begin{qoptions}
\opt{A}{H$_2$ \& D$_2$ have same internuclear distance}
\opt{B}{H$_2$ \& H$_3^+$ have same number of electrons}
\opt{C}{H$_3$O$^+$ \& H$_3^+$ have same number of 2c-2e$^-$ bonds}
\opt{D}{NaHCO$_3$(s) \& KHCO$_3$(s), have same arrangement of hydrogen bonded HCO$_3^-$ ions.}
\end{qoptions}
\qn{7} Identify correct order(s) of acidic strength?
\begin{qoptions}
\opt{A}{\qfigi{p1_chem_s1_q07_A}}
\opt{B}{\qfigi{p1_chem_s1_q07_B}}
\opt{C}{HF $<$ HCl $<$ HBr $<$ HI}
\opt{D}{HOCl $<$ HClO$_2$ $<$ HClO$_3$ $<$ HClO$_4$}
\end{qoptions}
\sect{SEC-2}
\paraq{8}{10}
{\small\noindent $\Delta H^0_{\text{atomisation}}[\mathrm{C_2H_6(g)}]=2810$ kJ/mol ;\\
$\Delta H^0_{\text{atomisation}}[\mathrm{C_2H_4(g)}]=2255$ kJ/mol\\
$\Delta H^0_f[\mathrm{C_2H_4(g)}]=25$ kJ/mol ;\\ $\Delta H^0_{\text{B.E.}}[\mathrm{H-H}]=430$ kJ/mol\\
Suppose : Bond enthalpy of same bond in diffrent molecule will be equal.}

\end{multicols}
\begin{center}\small\setlength{\tabcolsep}{4pt}
\begin{tabular}{|c|p{6cm}|c|p{3cm}|c|p{2.5cm}|}\hline
\multicolumn{2}{|c|}{Column-1} & \multicolumn{2}{c|}{Column-2} & \multicolumn{2}{c|}{Column-3}\\\hline
(A) & Enthalpy of hydrogenation of C$_2$H$_4$(g) & (P) & $\Delta H^0$ will be positive & (I) & $\Delta H^0>\Delta U^0$\\\hline
(B) & Enthalpy of formation of C$_2$H$_6$(g) & (Q) & $\Delta H^0=-225$ & (II) & $\Delta H^0<\Delta U^0$\\\hline
(C) & Enthalpy of hydrogenation of C$_6$H$_6$(g) (R.E. of C$_6$H$_6$(g) is $150$ kJ/mol) & (R) & $\Delta H^0=-125$ & (III) & $\Delta S^0<0$\\\hline
(D) & C$_4$H$_{10}$(g) $\to$ C$_2$H$_6$(g) $+$ C$_2$H$_4$(g) at very low temperature & (S) & $\Delta H^0=-100$ & (IV) & $\Delta G^0<0$\\\hline
\end{tabular}
\end{center}
\begin{multicols}{2}
\qn{1} Which of the following option correctly match
\begin{qoptions}
\opt{A}{A -- R -- II}
\opt{B}{A -- P -- I}
\opt{C}{A -- Q -- I}
\opt{D}{A -- S -- II}
\end{qoptions}
\qn{2} Which of the following option correctly match
\begin{qoptions}
\opt{A}{B -- P -- III}
\opt{B}{B -- Q -- III}
\opt{C}{B -- R -- III}
\opt{D}{B -- S -- III}
\end{qoptions}
\qn{3} Which of the following option incorrectly match
\begin{qoptions}
\opt{A}{D -- P -- I}
\opt{B}{D -- R -- I}
\opt{C}{C -- Q -- III}
\opt{D}{B -- S -- IV}
\end{qoptions}
\paraq{11}{13}
Answer Questions by appropriately matching the information given in the three columns of the following table.

\end{multicols}
\begin{center}\small\setlength{\tabcolsep}{4pt}
\begin{tabular}{|c|p{3cm}|c|p{4cm}|c|p{4.5cm}|}\hline
\multicolumn{2}{|c|}{\shortstack{Column - 1\\Geometry at\\bold atom O/N}} & \multicolumn{2}{c|}{\shortstack{Column - 2\\Observed on bold atom}} & \multicolumn{2}{c|}{\shortstack{Column - 3\\Also observed\\on bold atom}}\\\hline
(I) & Bent & (i) & H$_3$Si$\mathbf{N}$CO & (P) & P$_4\mathbf{O}_6$\\\hline
(II) & Pyramidal & (ii) & $\mathbf{O}_3$ & (Q) & {[}M(H$_2\mathbf{O})_6]^{+n}$\\\hline
(III) & Tetrahedral & (iii) & $\mathbf{N}$OF$_3$ & (R) & {[}$\mathbf{N}$O$_2]^+$\newline{[}NO$_3]^-$\\\hline
(IV) & Linear & (iv) & H$_3\mathbf{O}^+$ & (S) & {[}Be$_4\mathbf{O}$ (CH$_3$COO)$_6]$\\\hline
\end{tabular}
\end{center}
\begin{multicols}{2}
\qn{4} Select \textbf{correct} matching.
\begin{qoptions}
\opt{A}{(I) (ii) (R)}
\opt{B}{(II) (i) (R)}
\opt{C}{(III) (iii) (Q)}
\opt{D}{(IV) (i) (R)}
\end{qoptions}
\qn{5} Select \textbf{incorrect} matching.
\begin{qoptions}
\opt{A}{(III), (iii) (S)}
\opt{B}{(I), (ii) (P)}
\opt{C}{(II), (iv) (Q)}
\opt{D}{None}
\end{qoptions}
\qn{6} Select \textbf{correct} for anion of (R).
\begin{qoptions}
\opt{A}{Geometry is same as N of (i)}
\opt{B}{Hybridisation of same as N of (iii)}
\opt{C}{Bond angle is same as O of (ii)}
\opt{D}{Hybridisation is same as O of (ii)}
\end{qoptions}
\sect{SEC-3}
\qn{1} Total number of stereoisomers for following structure is $X$. Then identify the value of $(X-5)$ is :-
\qfig{p1_chem_s3_q01_1}
\qn{2} Identify total number of theoritically possible dichloro products formed in following reaction -
\qfig{p1_chem_s3_q02_1}
\qn{3} Find the number of paramagnetic compound(s) with octahedral anion.\\
K$_2$[CoF$_6$]; \ O$_2$[AsF$_6$]; \ K$_2$[CoCl$_4$]; \ K$_3$[Co(CN)$_6$];\\
K$_3$[CoF$_6$]; \ [Ni(NH$_3$)$_6$]SO$_4$; \ [Pt(NH$_3$)$_6$][Fe(CN)$_6$]; \ Na$_2[\mathrm{Fe^{II}(CN)_5NO}]$
\qn{4} Ionic solid AB(s) (M.wt. $=50$) crystallise in NaCl(s) type structure if edge length of unit cell is $\dfrac13\times10^{-7}$ cm then calculate density of AB(s) in gm/ml if it contain $\dfrac{100}{3}\%$ schottky defect [$N_A=6\times10^{23}$]
\qn{5} Identify reactions correctly matched with their major products -
\qfig{p1_chem_s3_q05_1}
\partsec{3}{MATHEMATICS}
\sect{SEC-1}
\qn{1} If $A$ is non-singular matrix of order $n$ such that $3ABA^{-1}+A=2A^{-1}BA$, then which of the following is/are true (where $P^{-1}$ denotes inverse of matrix $P$ and $|P|$ denotes determinant value of matrix $P$) -
\begin{qoptions}
\opt{A}{$|ABA^{-1}-A^{-1}BA|\neq0$}
\opt{B}{$|A+B|=0$}
\opt{C}{$|ABA^{-1}-A^{-1}BA|=0$}
\opt{D}{$|A+B|\neq0$}
\end{qoptions}
\qn{2} Normals are drawn from the point $P(15,12)$ to the parabola $y^2=4x$. If the feet of the normals form a $\triangle ABC$, then which of the following is/are true.
\begin{qoptions}
\opt{A}{Centroid of $\triangle ABC$ is $\left(\dfrac{26}3,\,0\right)$}
\opt{B}{Centroid of $\triangle ABC$ is $\left(\dfrac{22}3,\,0\right)$}
\opt{C}{Area of triangle formed by tangents to the parabola at $A$, $B$ and $C$ is $35$ sq. units}
\opt{D}{Area of triangle formed by tangents to the parabola at $A$, $B$ and $C$ is $140$ sq. units}
\end{qoptions}
\qn{3} Let $z_1$ lies on $|z|=1$ and $z_2$ lies on $|z|=2$, then which of the following is/are true -
\begin{qoptions}
\opt{A}{maximum value of $|z_1+z_2|$ is $3$}
\opt{B}{minimum value of $|2z_1-z_2|$ is $0$}
\opt{C}{maximum value of $|2z_1+z_2|$ is $4$}
\opt{D}{minimum value of $|2z_1-3z_2|$ is $5$}
\end{qoptions}
\qn{4} Two players $A$ and $B$ play a game which consist of a series of matches independent of each other. Whoever first win two matches, not necessarily consecutive, will win the game. The proability of $A$'s winning, drawing or losing a match against $B$ are $\dfrac12,\dfrac13,\dfrac16$ respectively. Which of the following is/are true -
\begin{qoptions}
\opt{A}{Probability of $A$ winning the game at the end of $11^{\text{th}}$ match is
\[{}^{10}C_1\left(\frac12\right)^2\left(\frac13\right)^9+{}^{10}C_2\left(\frac12\right)^2\left(\frac13\right)^9\]}
\opt{B}{Probability of $A$ winning the game at the end of $11^{\text{th}}$ match is
\[{}^{10}C_1\left(\frac12\right)^2\left(\frac13\right)^9+{}^{10}C_2\left(\frac12\right)\left(\frac13\right)^9\]}
\opt{C}{If it is known that $A$ win the game at the end of $11^{\text{th}}$ match then the probability that $B$ will win only one match is $\dfrac9{11}$}
\opt{D}{If it is known that $A$ win the game at the end of $11^{\text{th}}$ match then the probability that $B$ will win only one match is $\dfrac9{13}$}
\end{qoptions}
\qn{5} Let $m$ is the number of points where $f(x)=\{x\}+\{2x\}+\{3x\}$ is non-differentiable in $[0,n)$ (where $m,n\in N$ and $\{.\}$ denotes fractional part function), then which of the following is/are true -
\begin{qoptions}
\opt{A}{If $n=1$, then $m=3$}
\opt{B}{If $n=3$, then $m=11$}
\opt{C}{If $n=2$, then $m=7$}
\opt{D}{If $n=3$, then $m=12$}
\end{qoptions}
\qn{6} If pair of tangents are drawn from centre $C$ of the circle $x^2+y^2-5x=0$ to the nearer branch of the hyperbola $\dfrac{x^2}{25}-\dfrac{y^2}{16}=1$ to touch it at point $A$ and $B$, then
\begin{qoptions}
\opt{A}{product of the slopes of lines $CA$ and $CB$ is $\dfrac{-64}{75}$}
\opt{B}{product of the slopes of lines $CA$ and $CB$ is $\dfrac{-32}{25}$}
\opt{C}{circumcentre of $\triangle ABC$ lies inside the $\triangle ABC$.}
\opt{D}{circumcentre of $\triangle ABC$ lies outside the $\triangle ABC$.}
\end{qoptions}
\qn{7} If $a+b+c=4$ and $a^2+b^2+c^2+3(ab+bc+ac)=21$ where $a,b,c\in R$, then which of the following is/are true
\begin{qoptions}
\opt{A}{minimum value of $(a+b)(b+c)(a+c)$ is $18$}
\opt{B}{minimum value of $(a+b)(b+c)(a+c)$ is $24$}
\opt{C}{maximum value of $(a+b)(b+c)(a+c)$ is $\dfrac{490}{27}$}
\opt{D}{maximum value of $(a+b)(b+c)(a+c)$ is $\dfrac{471}{27}$}
\end{qoptions}
\sect{SEC-2}
\paraq{8}{10}
Consider $f(x)=xe^{-x}$, $g(x)=xe^{x}$, $p(x)=x^3-3x$ and $q(x)=x^4-2x^2+3$

Column-1 contain information about different functions.

Column-2 and Column-3 contain information about various aspects of $h(x)$.

\end{multicols}
\begin{center}\small\setlength{\tabcolsep}{4pt}
\begin{tabular}{|c|p{3cm}|c|p{4.5cm}|c|p{5cm}|}\hline
\multicolumn{2}{|c|}{Column-1} & \multicolumn{2}{c|}{Column-2} & \multicolumn{2}{c|}{Column-3}\\\hline
(I) & $h(x)=f(x)$ & (i) & $h(x)$ is increasing in $x\in(0,1)$ & (P) & $h(x)=2$ has $2$ solution.\\\hline
(II) & $h(x)=g(x)$ & (ii) & $h(x)$ is decreasing in $x\in(0,1)$ & (Q) & $h(x)=0$ has $3$ solutions.\\\hline
(III) & $h(x)=p(x)$ & (iii) & $h'(x)$ is increasing in $x\in(-\infty,-2)$ & (R) & $h''(x)<0$ for $x\in(-\infty,-3)$.\\\hline
(IV) & $h(x)=q(x)$ & (iv) & $h'(x)$ is decreasing in $x\in(-\infty,-2)$ & (S) & $h(x)$ has a negative point of minima.\\\hline
\end{tabular}
\end{center}
\begin{multicols}{2}
\qn{1} Which of the following option is the only \textbf{CORRECT} combination ?
\begin{qoptions}
\opt{A}{(I) (iv) (Q)}
\opt{B}{(IV) (iii) (R)}
\opt{C}{(II) (iv) (Q)}
\opt{D}{(I) (iv) (R)}
\end{qoptions}
\qn{2} Which of the following option is the only \textbf{INCORRECT} combination ?
\begin{qoptions}
\opt{A}{(II) (i) (R)}
\opt{B}{(III) (ii) (P)}
\opt{C}{(II) (ii) (R)}
\opt{D}{(III) (iv) (R)}
\end{qoptions}
\qn{3} Which of the following option is the only \textbf{CORRECT} combination ?
\begin{qoptions}
\opt{A}{(IV) (i) (S)}
\opt{B}{(IV) (iii) (P)}
\opt{C}{(II) (iv) (Q)}
\opt{D}{(III) (ii) (S)}
\end{qoptions}
\paraq{11}{13}
Column-1 contain the hyperbola with the given conditions.

Column-2 contain the eccentricity of the hyperbola.

Column-3 contain the information about tangents to the hyperbola.

\end{multicols}
\begin{center}\small\setlength{\tabcolsep}{4pt}
\begin{tabular}{|c|p{5.5cm}|c|p{1.2cm}|c|p{5.5cm}|}\hline
\multicolumn{2}{|c|}{Column 1} & \multicolumn{2}{c|}{Column 2} & \multicolumn{2}{c|}{Column 3}\\\hline
(I) & The distance of one focus of the hyperbola $\dfrac{x^2}{a^2}-\dfrac{y^2}{b^2}=1$ from its directrix is $5$ and $3$ & (i) & $\sqrt{17}$ & (P) & There exist infinite points on the plane from which perpendicular tangents can be drawn to hyperbola.\\\hline
(II) & Equation of hyperbola is $x=\dfrac{e^t+e^{-t}}{2}$ and $y=\dfrac{e^t-e^{-t}}{3}$, $t\in R$. & (ii) & $\sqrt{\dfrac32}$ & (Q) & There exist no point on the plane from which perpendicular tangents can be drawn to hyperbola.\\\hline
(III) & Tangent to hyperbola $\dfrac{x^2}{a^2}-\dfrac{y^2}{b^2}=1$ at point $\left(\dfrac{2a}{\sqrt3},\dfrac{b}{\sqrt3}\right)$ form a triangle of area $3a^2$ square units with the co-ordinate axes. & (iii) & $\dfrac{\sqrt{13}}3$ & (R) & There exist a tangent to hyperbola having slope $1$.\\\hline
(IV) & The normal at the point $(6,3)$ on the hyperbola $\dfrac{x^2}{a^2}-\dfrac{y^2}{b^2}=1$ intersect the $x$-axis at $(9,0)$. & (iv) & $2$ & (S) & There exist a tangent to hyperbola having slope $2$.\\\hline
\end{tabular}
\end{center}
\begin{multicols}{2}
\qn{4} Which of the following option is the only \textbf{CORRECT} combination ?
\begin{qoptions}
\opt{A}{(I) (iv) (R)}
\opt{B}{(II) (iii) (Q)}
\opt{C}{(III) (i) (Q)}
\opt{D}{(IV) (ii) (Q)}
\end{qoptions}
\qn{5} Which of the following option is the only \textbf{INCORRECT} combination ?
\begin{qoptions}
\opt{A}{(I) (iv) (P)}
\opt{B}{(II) (iii) (P)}
\opt{C}{(IV) (ii) (R)}
\opt{D}{(III) (i) (Q)}
\end{qoptions}
\qn{6} Which of the following option is the only \textbf{INCORRECT} combination ?
\begin{qoptions}
\opt{A}{(I) (iv) (S)}
\opt{B}{(II) (iii) (S)}
\opt{C}{(IV) (ii) (P)}
\opt{D}{(III) (i) (S)}
\end{qoptions}
\sect{SEC-3}
\qn{1} Consider the number $N=10!$. If the number of positive divisors of $N$ is $\lambda$, then $\dfrac{\lambda}{30}$ equals
\qn{2} Let $L_1$, $L_2$ and $L_3$ be the length of tangents drawn from a point $P$ to the circles $x^2+y^2=4$, $x^2+y^2-4x=0$ and $x^2+y^2-4y=0$ respectively. If $L_1^4=L_2^2L_3^2+16$, then locus of $P$ are the curves $C_1$ (a straight line) and $C_2$ (a circle). Let the circumcentre of the triangle formed by $C_1$ and other two lines which are inclined at an angle of $45^\circ$ with $C_1$ and is tangent to $C_2$ is $(\alpha,\beta)$, then $(\alpha+2\beta)$ equals
\qn{3} Given $a,b\in\{0,1,2,\dots,6\}$. Consider the system of equations
\begin{gather*}
x+y+z=4,\\
2x+y+3z=6,\\
x+2y+az=b
\end{gather*}
If the number of ordered pairs $(a,b)$ so that the system of equations has unique solution is $\lambda$, then $\dfrac\lambda7$ equals
\qn{4} Let the function $g(x)$ where $g:R\to R$, continuous in $[0,\infty)$ and positive in $(0,\infty)$ satisfying the condition $g(1)=1$ and $\dfrac12\displaystyle\int_0^x g^2(t)\,dt=\dfrac1x\left(\displaystyle\int_0^x g(t)\,dt\right)^2$ then the value of $\left(g(4)\right)^{\sqrt2-1}$ is
\qn{5} Let $S_n=\displaystyle\sum_{r=1}^{n}r\cdot2^{\frac r2(1+(-1)^r)}\cdot3^{\frac r2(1-(-1)^r)}$. If $S_{20}=\alpha\cdot3^{21}+\beta\cdot2^{22}+\dfrac{391}{288}$, then the value of $\left[\dfrac{32\alpha-9\beta}{4}\right]$ is\\
(where $[.]$ denotes the greatest integer function)
\papersec{2}
\gdef\paraq#1#2{\par\smallskip\noindent\textbf{Paragraph for Questions #1 and #2}\par}
\partsec{1}{PHYSICS}
\sect{SEC-1}
\qn{1} A uniform pole is attached to a vertical wall by a frictionless pivot. The pole is held horizontal by a vertical massless cable attached to the ceiling as shown in the diagram. Considering torques on the pole about the axis of the pivot, which of the statement is \textbf{CORRECT} ?
\qfig{p2_phy_s1_q01_1}
\begin{qoptions}
\opt{A}{The magnitude of the torque due to tension in the cable is greater than the magnitude of torque due to weight of pole.}
\opt{B}{The magnitude of the torque due to pivot on the pole is equal to the magnitude of torque due to weight of pole.}
\opt{C}{The magnitude of torque due to tension in cable is equal to magnitude of torque due weight of pole.}
\opt{D}{Tension in cable is more than the weight of the pole}
\end{qoptions}
\qn{2} For progressive wave : $y=(2\,\text{cm})\sin(3t-4x)$\\
\textbf{Statement-1 :} The particle at $x=0$ is moving towards positive $y$-axis at $t=0$.\\
\textbf{Statement-2 :} The particle at $x=8$m is in phase with particle at $x=1$m.\\
\textbf{Statement-3 :} The acceleration of particle at $x=\dfrac\pi8$m at $t=\dfrac\pi6$ sec is increasing in magnitude.\\
\textbf{Statement-4 :} The velocity of particle at $x=\dfrac\pi6$m is in direction of propagation of wave at all times.\\
The correct statement are :
\begin{qoptions}
\opt{A}{$1$}
\opt{B}{$1$ and $3$}
\opt{C}{$1$, $2$ and $3$}
\opt{D}{$1$, $2$, $3$ and $4$}
\end{qoptions}
\qn{3} An organ pipe open at both ends is oscillating in fifth overtone. The separation between two points where pressure amplitude is $\dfrac1{\sqrt2}$ times the pressure amplitude at the anti-nodes is $10$cm. The length of the organ tube is : (Take speed of sound in air to be $330$ m/s)
\begin{qoptions}
\opt{A}{$40$ cm}
\opt{B}{$80$ cm}
\opt{C}{$100$ cm}
\opt{D}{$120$ cm}
\end{qoptions}
\qn{4} A wire of uniform resistance per unit length is bent to form an equilateral triangle of side $L$. A current $I$ flows into one corner and flows out of an adjacent corner, as shown. The magnetic field at the centroid of triangle due to the current in the triangular frame is
\qfig{p2_phy_s1_q04_1}
\begin{qoptions}
\opt{A}{$3\mu_0I/2\pi L$}
\opt{B}{$3\mu_0I/\pi L$}
\opt{C}{$3\sqrt3\mu_0I/2\pi L$}
\opt{D}{zero}
\end{qoptions}
\qn{5} Rank the vernier callipers/screw gauge in the order of their increasing precision :\\
(1) Screw gauge having main scale division $1$mm and no. of circular division $=200$.\\
(2) Screw gauge having main scale division $0.5$ mm and no. of circular division $=400$.\\
(3) Vernier callipers having main scale division $=1$mm such that $9$ MSD $=10$ VSD.\\
(4) Vernier callipers having main scale division $=1$mm such that $4$ MSD $=5$ VSD.
\begin{qoptions}
\opt{A}{$4<2<3<1$}
\opt{B}{$2<4<1<3$}
\opt{C}{$4<3<2<1$}
\opt{D}{$4<3<1<2$}
\end{qoptions}
\qn{6} A YDSE experiment is used to measure the wavelength of light used. The distance between screen and slits is measured by a meter scale with least count $1$mm. The measured distance is $1.500$ m. The distance between the slits is measured by a micrometer of least count $0.01$ mm. The distance between slits is measured to be $0.33$ mm. The fringe width is measured by the same micrometer. It comes out to be $2.50$ mm. The wavelength of light is :-
\begin{qoptions}
\opt{A}{$550\pm19$ nm}
\opt{B}{$550\pm39$ nm}
\opt{C}{$330\pm28$ nm}
\opt{D}{$330\pm56$ nm}
\end{qoptions}
\qn{7} Young's double slit experiment is conducted in a liquid of refractive index $\mu_1$ as shown in figure. A thin transprent slab of refractive index $\mu_2$ and thickness $t$ is placed in front of the slit $s_2$. The magnitude of optical path difference at $O$ is :-
\qfig{p2_phy_s1_q07_1}
\begin{qoptions}
\opt{A}{$\left|\left(\dfrac{\mu_2}{\mu_1}-1\right)t\right|$}
\opt{B}{$\left|\left(\dfrac{\mu_1}{\mu_2}-1\right)t\right|$}
\opt{C}{$\left|(\mu_2-1)\,t\right|$}
\opt{D}{$\left|(\mu_2-\mu_1)\,t\right|$}
\end{qoptions}
\sect{SEC-2}
\qn{1} A battery of emf $\varepsilon_0=5$V and internal resistance $5\Omega$ is connected across a long uniform wire $AB$ of length $1$ m and resistance per unit length $5\ \Omega\text{m}^{-1}$. Two cells of $\varepsilon_1=1$V and $\varepsilon_2=2$V are connected as shown in the figure.
\qfig{p2_phy_s2_q01_1}
\begin{qoptions}
\opt{A}{The null point is at $A$}
\opt{B}{If the Jockey is touched to point $B$ the current in the galvanometer will be going towards $B$}
\opt{C}{When Jockey is connected to point $A$ no current is flowing through $1$ V battery}
\opt{D}{The null point is at distance of $8/15$ m from $A$.}
\end{qoptions}
\qn{2} A loop $PQR$ formed by three identical uniform conducting rods each of length $a$ is suspended from one of its vertices ($P$) so that it can rotate about horizontal fixed smooth axis $CD$. Initially plane of loop is in vertical plane. A constant current $i$ is flowing in the loop. Total mass of the loop is $m$. At $t=0$, a uniform magnetic field of strength $B$ directed vertically upwards is switched on. Acceleration due to gravity is $g$.
\qfig{p2_phy_s2_q02_1}
\begin{qoptions}
\opt{A}{Minimum value of $B$ so that the plane of the loop becomes horizontal (even for an instant) during its subsequent motion is $\dfrac{2mg}{3ia}$.}
\opt{B}{Minimum value of $B$ so that the plane of the loop becomes horizontal (even for an instant) during its subsequent motion is $\dfrac{4mg}{3ia}$.}
\opt{C}{The equilibrium position (angle of the plane of the loop with vertical) is given by $\tan^{-1}\left(\dfrac{3Bai}{2mg}\right)$.}
\opt{D}{The equilibrium position (angle of the plane of the loop with vertical) is given by $\tan^{-1}\left(\dfrac{3Bai}{4mg}\right)$.}
\end{qoptions}
\qn{3} A spherical shell of radius $R$ is given charge $3q$ on its surface and a point charge $q$ is placed at distance $R/2$ from its centre $C$. Also there is a charge $2q$ placed outside the shell at a distance of $2R$ as shown. Then :
\qfig{p2_phy_s2_q03_1}
\begin{qoptions}
\opt{A}{The magnitude of electric field at the centre $C$ due to charges on the outer surface of shell is $\dfrac{Kq}{2R^2}$ before closing the switch $S$.}
\opt{B}{The electric potential at the centre $C$ due to charges on the outer surface of shell is $\left(\dfrac{-Kq}{R}\right)$ before closing the switch $S$.}
\opt{C}{The electric potential at the centre $C$ due to charges on the outer surface of shell is $\left(\dfrac{-Kq}{R}\right)$ after closing the switch $S$.}
\opt{D}{Charge flown through the switch into earth after closing the switch $S$ is $5q$.}
\end{qoptions}
\qn{4} Three small balls $A$, $B$ and $C$ of same mass $m$ are fixed by massless rods at the vertices of equilateral triangle with side $L$. The system is placed on a smooth horizontal surface and set in rotation about the stationary COM with angular speed $\omega$. At a certain instant ball $B$ breaks off the system with a velocity $\vec v_B$ as shown in diagram where $|\vec v_B|=\dfrac{2\omega L}{\sqrt3}$.
\qfig{p2_phy_s2_q04_1}
\begin{qoptions}
\opt{A}{The COM of $A$ and $C$ moves with a velocity $\dfrac{-\vec v_B}{2}$.}
\opt{B}{The tension in rod connecting $A$ and $C$ after $B$ breaks off is $m\omega^2\dfrac\ell2$}
\opt{C}{The speed of $A$ just after $B$ breaks off is $\dfrac{v_B}{2}$}
\opt{D}{The speed of $A$, after a time $\dfrac\pi\omega$ after $B$ breaks off is $\dfrac{\sqrt{19}}{6}\omega\ell$}
\end{qoptions}
\qn{5} A syringe is filled with water. Its volume is $20$ cm$^3$ and cross-section of its interior part is $4$cm$^2$. The syringe is held vertically such that its nozzle is at its top and its $100$g piston is pressed by external agent and moves it with a constant speed. The ejected water has a initial upward velocity of $2$m/s and cross-section of the beam of water at the nozzle is $1$mm$^2$.
\begin{qoptions}
\opt{A}{The speed with which the piston is moving is $5$ mm/sec.}
\opt{B}{The external force acting on the piston is $42$ N.}
\opt{C}{The total work done by external agent in emptying the syringe is $0.095$ J}
\opt{D}{The external force remains constant during the process.}
\end{qoptions}
\qn{6} In the circuit shown switch $S$ is closed at $t=0$. At time ($t$) which of the following is/are correct:
\qfig{p2_phy_s2_q06_1}
\begin{qoptions}
\opt{A}{$I=\dfrac V{2R}\left(1-e^{-\frac{4Rt}{L}}\right)$}
\opt{B}{$I_1=I_2=\dfrac V{2R}\left(1-e^{-\frac{4Rt}{L}}\right)$}
\opt{C}{At $t=0$, $\dfrac{dV_{ab}}{dt}=0$}
\opt{D}{$I=\sqrt{I_1^2+I_2^2}$}
\end{qoptions}
\qn{7} H like atom of atomic number $Z$ is in an excited state of quantum number $2n$. It can emit a maximum energy photon of $204$ eV. If it makes a transition to quantum state $n$, a photon of energy $40.8$ eV is emitted. Identify the correct statement/s
\begin{qoptions}
\opt{A}{$n=2$, $Z=16$}
\opt{B}{$n=2$, $Z=4$}
\opt{C}{ground state energy for this atom is $-3481.6$ eV}
\opt{D}{ground state energy for this atom is $-217.6$ eV}
\end{qoptions}
\sect{SEC-3}
\paraq{15}{16}
One mole of an ideal monoatomic gas is taken along a cyclic process : (take : $\pi=3.14$)
\qfig{p2_phy_s3_para15_16_1}
Process (1 - 2) : One fourth of the circle\\
Process (2 - 3) : Isochoric\\
Process (3 - 1) : Isobaric
\qn{1} Net work done around the cyclic process is :
\begin{qoptions}
\opt{A}{$P_0V_0$}
\opt{B}{$\dfrac{P_0V_0}3$}
\opt{C}{$0.22\,P_0V_0$}
\opt{D}{$\dfrac45P_0V_0$}
\end{qoptions}
\qn{2} The efficiency of cyclic process is :
\begin{qoptions}
\opt{A}{$0.4$}
\opt{B}{$0.08$}
\opt{C}{$0.04$}
\opt{D}{$0.57$}
\end{qoptions}
\paraq{17}{18}
A ray of light enters a spherical drop of water of index $n$ as shown in figure. Light is partially reflected from the rear surface, the reflected light re-enters air after refraction as shown. The angle between incoming ray and outgoing ray is known as angle of deflection. ($n<1$)
\qfig{p2_phy_s3_para17_18_1}
\qn{3} Expression for the angle of deflection $\delta$ (in clockwise) is
\begin{qoptions}
\opt{A}{$\pi-4\alpha-2\phi$}
\opt{B}{$\pi-4\alpha+2\phi$}
\opt{C}{$2\phi$}
\opt{D}{$4\alpha$}
\end{qoptions}
\qn{4} The angle $\phi$ which produces minimum deflection is :-
\begin{qoptions}
\opt{A}{$\cos^{-1}\sqrt{\dfrac{n^2}{3}}$}
\opt{B}{$\cos^{-1}\sqrt{\dfrac{n}{3}}$}
\opt{C}{$\cos^{-1}\sqrt{\dfrac{n^2-1}{3}}$}
\opt{D}{$\cos^{-1}\sqrt{\dfrac{n^2-2}{3}}$}
\end{qoptions}
\partsec{2}{CHEMISTRY}
\sect{SEC-1}
\qn{1} Choose the major product of the following reaction sequence
\qfig{p2_chem_s1_q01_1}
\begin{qoptions}
\opt{A}{\qfigi{p2_chem_s1_q01_A}}
\opt{B}{\qfigi{p2_chem_s1_q01_B}}
\opt{C}{\qfigi{p2_chem_s1_q01_C}}
\opt{D}{\qfigi{p2_chem_s1_q01_D}}
\end{qoptions}
\qn{2} A hexapeptide has the composition Ala, Gly, Phe, Val$_3$. Both the N-terminal and C-terminal units are Val. Cleavage of the hexapeptide by chymotrypsin gives two different tripeptides, both having Val as the N-terminal group (Chymotrypsin hydrolyzes on the C-side of amino acids that contain aromatic six membered ring). Among the products of random hydrolysis is a Ala-Val dipeptide fragment. What is the primary structure of the hexapeptide?
\begin{qoptions}
\opt{A}{Val-Gly-Phe-Val-Ala-Val}
\opt{B}{Val-Ala-Phe-Val-Gly-Val}
\opt{C}{Val-Gly-Ala-Val-Phe-Val}
\opt{D}{Val-Phe-Val-Ala-Gly-Val}
\end{qoptions}
\qfullbegin
\qn{3}
\qfig{p2_chem_s1_q03_1}
\qfullsplit
\textbf{INCORRECT} geometrical arrangement of reactant complex is.
\begin{qoptions}
\opt{A}{\qfigi{p2_chem_s1_q03_A}}
\opt{B}{\qfigi{p2_chem_s1_q03_B}}
\opt{C}{\qfigi{p2_chem_s1_q03_C}}
\opt{D}{\qfigi{p2_chem_s1_q03_D}}
\end{qoptions}
\qfullend
\qn{4} Select incorrect for given reaction sequence.
\begin{gather*}
\mathrm{P_4+I_2+H_2O\to(X)+HI+H_3PO_4}\\
\mathrm{(X)+KOH_{aq.}\to(Y)_{(g)}+Z_{aq.}+H_2O}
\end{gather*}
$\mathrm{(Y)_{(g)}}\xrightarrow{\mathrm{HgCl_2}}$ ppt\\
$\mathrm{(Y)_{(g)}}\xrightarrow{\mathrm{CuSO_{4\,aq.}}\,(1\text{-}2\text{ drop})}$ white ppt $+$ brown colour solution
\begin{qoptions}
\opt{A}{$\mathrm{Y_{(g)}}$ produce ppt with CuSO$_4$}
\opt{B}{Reaction of $\mathrm{Z_{aq.}}$ + CuSO$_4$ is redox change}
\opt{C}{Bond angle in P$_4$ $<$ $\mathrm{(Y)_{(g)}}$}
\opt{D}{On heating H$_3$PO$_4$ it produce $\mathrm{Y_{(g)}}$}
\end{qoptions}
\qn{5} Given hydrogen-oxygen fuel cell -
\qfig{p2_chem_s1_q05_1}
\begin{gather*}
\mathrm{H_2(g)}+\tfrac12\mathrm{O_2(g)}\to\mathrm{H_2O(\ell)};\ E^0_{\text{cell}}=1.23\,\text{V}\\
(\Delta_fH)_{\mathrm{H_2O(g)}}=-245\ \text{kJmol}^{-1},\\
\Delta_{\text{vap}}H(\mathrm{H_2O},\ell)=40\ \text{kJmol}^{-1}\\
\log2=0.3,\ \frac{2.303RT}{F}=0.06
\end{gather*}
The efficiency of fuel cell is -
\begin{qoptions}
\opt{A}{$96.08$ \%}
\opt{B}{$41.3$ \%}
\opt{C}{$61.95$\%}
\opt{D}{$82.68$\%}
\end{qoptions}
\qn{6} Which is incorrect statement ?
\begin{qoptions}
\opt{A}{Brownian movement is responsible for the stability of colloidal solution}
\opt{B}{No Tyndall effect observed in True solution}
\opt{C}{Haemoglobin (blood) and gold sol are negative charge colloids}
\opt{D}{Casein act as emulsifier for milk}
\end{qoptions}
\qn{7} Given :\\
\begin{gather*}\mathrm{H_2O_2+KMnO_4+H_2SO_4\to}\\\mathrm{MnSO_4+H_2O+O_2+K_2SO_4}\end{gather*}
which is correct statement ?
\begin{qoptions}
\opt{A}{Number of gm equivalent of H$_2$O$_2$ and KMnO$_4$ reacts in the ratio of $5:2$}
\opt{B}{Number of moles of MnSO$_4$ and O$_2$ produced in the ratio of $5:2$}
\opt{C}{Equivalent mass of H$_2$SO$_4$ is $29.4$}
\opt{D}{$100$ ml $2.27$ V H$_2$O$_2$ reacts completely with $160$ ml $0.5$M KMnO$_4$ in given reaction}
\end{qoptions}
\sect{SEC-2}
\qn{1} Which of following gives same product by S$_\mathrm{N}1$ and S$_\mathrm{N}2$ (excluding stereo)
\begin{qoptions}
\opt{A}{\qfigi{p2_chem_s2_q01_A}}
\opt{B}{\qfigi{p2_chem_s2_q01_B}}
\opt{C}{\qfigi{p2_chem_s2_q01_C}}
\opt{D}{\qfigi{p2_chem_s2_q01_D}}
\end{qoptions}
\qn{2} Select the correct statement -
\begin{qoptions}
\opt{A}{Area enclosed by a cycle on $P$-$V$ diagram is same as on $T$-$S$ diagram}
\opt{B}{$\Delta S=nC_{pm}\ln\dfrac{T_2}{T_1}$ can be used in isobaric process for both ideal and real gas}
\opt{C}{$\Delta S=nC_{pm}\ln\dfrac{T_2}{T_1}+nR\ln\dfrac{P_1}{P_2}$ is applicable for both reversible and irreversible process for an ideal gas not undergoing any chemical and phase change.}
\opt{D}{Work done on a diatomic gas will be more than monoatomic gas compressed reversible adiabatically from same initial state to same final volume.}
\end{qoptions}
\qn{3} As per Fajan's rule polarization (covalent character) will be increased by?
\begin{qoptions}
\opt{A}{High charge and small size cation}
\opt{B}{High charge and large size anion}
\opt{C}{Low charge and small size anion}
\opt{D}{Low charge and large size cation}
\end{qoptions}
\qn{4} The following first order reaction -\\
$2\mathrm{A(aq.)}\to3\mathrm{B(aq.)}+4\mathrm{C(aq.)}$\\
is monitored by measuring optical rotation of reaction mixture containing optical active dextrorotatory inert impurity at different time interval. The species A, B and C are optically active with specific rotation $20^\circ$, $40^\circ$ and $-80^\circ$ per mol respectively.

\begin{center}\small
\begin{tabular}{@{}lccc@{}}
Time (min) & $0$ & $22$ & $\infty$\\
Optical rotation of reaction mixture & $120^\circ$ & $-40^\circ$ & $-120^\circ$
\end{tabular}
\end{center}
(Given - $\ln3=1.1$, $\ln2=0.7$)\\
which is/are correct statements -
\begin{qoptions}
\opt{A}{Rate constant of reaction is $0.05$ min$^{-1}$}
\opt{B}{Rate constant of apperance of C is $0.1$ min$^{-1}$}
\opt{C}{The time in which $75\%$ of A reacted is $28$ min}
\opt{D}{The time in which $50\%$ of A reacted is $420$ sec.}
\end{qoptions}
\qn{5} Reagent which can't be used to carryout this conversion are :
\qfig{p2_chem_s2_q05_1}
\begin{qoptions}
\opt{A}{LiAlH$_4$}
\opt{B}{N$_2$H$_4$/OH$^-$/$\Delta$}
\opt{C}{N$_2$H$_4$/H$_2$O$_2$}
\opt{D}{Zn(Hg)HCl}
\end{qoptions}
\qn{6} Select correct statement(s).
\begin{qoptions}
\opt{A}{Blue solution of Na in liq. NH$_3$ is unstable with respect to amide formation}
\opt{B}{Na$_2$O$_2$ react with CO and evolve O$_2$}
\opt{C}{Li$^+$ does not form alum due to its high hydration energy}
\opt{D}{NaNO$_3$ when heated at $800^\circ$C produce two gases both are paramagnetic}
\end{qoptions}
\qn{7} Given :
\begin{gather*}\Psi_{\text{radial}}:\ \frac1{9\sqrt2}\left(\frac1{a_0}\right)^{3/2}\\\times(\sigma^2-4\sigma+3)\,e^{-\sigma/2}\end{gather*}
where $\sigma=\dfrac{2r}{a_0}$, $a_0=0.529$\AA\\
which is/are correct statement -
\begin{qoptions}
\opt{A}{\qfigi{p2_chem_s2_q07_A}}
\opt{B}{\qfigi{p2_chem_s2_q07_B}}
\opt{C}{\qfigi{p2_chem_s2_q07_C}}
\opt{D}{Distance between $2^{\text{nd}}$ radial nodes and $1^{\text{st}}$ radial nodes is equal to $0.529$\AA}
\end{qoptions}
\sect{SEC-3}
\paraq{15}{16}
Metal ($\mathbf M$) is a white, crystalline metal and is not appreciably ductile or malleable. It has very high melting point. The metal is soluble in dilute or concentrated hydrochloric acid. If air is excluded, Metal(II) ions are formed :
\begin{gather*}
\mathrm{M+2H^+\to M^{2+}+H_2\uparrow}\\
\mathrm{M+2HCl\to M^{2+}+2Cl^-+H_2\uparrow}
\end{gather*}
Metal(II) ions are rather unstable as they are strong reducing agent, they decompose even in acidic solution with the formation of hydrogen and $\mathrm{M^{+3}}$.

In the presence of atomospheric oxygen, metal ($\mathbf M$) gets partly or wholly oxidized to the tervalent state :
\[\mathrm{4M^{2+}+O_2+4H^+\to4M^{3+}+2H_2O}\]
Dilute sulphuric acid attacks metal ($\mathbf M$) slowly, with the formation of hydrogen. In hot, concentrated sulphuric acid metal ($\mathbf M$) dissolves readily, to form metal(III) ions and sulphur dioxide:
\[\mathrm{2M+6H_2SO_4\to2M^{3+}+3SO_4^{2-}+3SO_2\uparrow+6H_2O}\]
Both dilute and concentrated nitric acid render metal ($\mathbf M$) passive.
\qn{1} Colour of $\mathrm{M^{+2}}$ ion \& $\mathrm{M^{+3}}$ ion respectively when obtain by dissolving MO \& M$_2$O$_3$ in dil. HCl.
\begin{qoptions}
\opt{A}{Blue, Green}
\opt{B}{Green, Blue}
\opt{C}{Blue, Yellow}
\opt{D}{Violet, Blue}
\end{qoptions}
\qn{2} Aqueous solution of MCl$_3$.6H$_2$O when reacts with Na$_2$O$_2$ it produce coloured solution(s). Coloured solution(s) does not produce any ppt with (BaCl$_2$ + dil. HCl), but yellow ppt is produced with (BaCl$_2$ + CH$_3$COOH$_{\text{aq.}}$).\\
Select incorrect about $\mathbf M$Cl$_3$.6H$_2$O.
\begin{qoptions}
\opt{A}{Produce violet pink appearance with excess NH$_4$OH in aq.}
\opt{B}{Produce light green solution with excess NaOH$_{\text{aq.}}$}
\opt{C}{Produce M$_2$S$_3$ when react with aq. solution of (NH$_4$)$_2$S}
\opt{D}{It is paramagnetic}
\end{qoptions}
\paraq{17}{18}
One of the best methods for the regioselective alkylation of ketones consists in an activation of the required $\alpha$-carbon atom. For example :
\qfig{p2_chem_s3_para17_18_1}
\qn{3} Identify structure of $\mathrm{D_2}$?
\begin{qoptions}
\opt{A}{\qfigi{p2_chem_s3_q03_A}}
\opt{B}{\qfigi{p2_chem_s3_q03_B}}
\opt{C}{\qfigi{p2_chem_s3_q03_C}}
\opt{D}{\qfigi{p2_chem_s3_q03_D}}
\end{qoptions}
\qn{4} Identify structure of $\mathrm{D_1}$?
\begin{qoptions}
\opt{A}{\qfigi{p2_chem_s3_q04_A}}
\opt{B}{\qfigi{p2_chem_s3_q04_B}}
\opt{C}{\qfigi{p2_chem_s3_q04_C}}
\opt{D}{\qfigi{p2_chem_s3_q04_D}}
\end{qoptions}
\partsec{3}{MATHEMATICS}
\sect{SEC-1}
\qn{1} If $E_1$ and $E_2$ are two events such that $P(E_1)=\dfrac14$, $P\!\left(\dfrac{E_2}{E_1}\right)=\dfrac12$ and $P\!\left(\dfrac{E_1}{E_2}\right)=\dfrac14$, then which one is \textbf{INCORRECT}
\begin{qoptions}
\opt{A}{$E_1$ and $E_2$ are mutually independent events.}
\opt{B}{$E_1$ and $E_2$ are exhaustive events}
\opt{C}{$E_2$ is twice as likely to occur as $E_1$}
\opt{D}{$P(E_1\cap E_2)$, $P(E_1)$, $P(E_2)$ are in G.P.}
\end{qoptions}
\qn{2} Let $N$ be the number of $5$-digit nutural numbers beginning with $9$ that has exactly three identical digits (other digits are distinct). Then the sum of the digits of $N$ is equal to -
\begin{qoptions}
\opt{A}{$9$}
\opt{B}{$8$}
\opt{C}{$7$}
\opt{D}{$6$}
\end{qoptions}
\qn{3} Let $f(x)$ be a polynomial function of degree five satisfying\\
$f(|x|)=|f(x)|$, $\forall\ x\in R$\\
$f(1)=0$, $f(2)=18$\\
then which of the following is NOT correct
\begin{qoptions}
\opt{A}{$x=1$ is a twice repeated root of $f(x)=0$}
\opt{B}{$x=-1$ is a twice repeated root of $f(x)=0$}
\opt{C}{The equation $f(x).f''(x)+(f'(x))^2=0$ can have less than $4$ distinct real roots.}
\opt{D}{The equation $f(x).f''(x)+(f'(x))^2=0$ can not have less than $4$ distinct real roots.}
\end{qoptions}
\qn{4} If number of $5$-digit numbers of the form $abcde$ where $a,b,c,d,e\in\{0,1,2,\dots,9\}$ and $b=a+c$, $d=c+e$ is $\lambda$, then $\left[\dfrac\lambda{100}\right]$ is (where $[.]$ represents greatest integral function)
\begin{qoptions}
\opt{A}{$4$}
\opt{B}{$5$}
\opt{C}{$3$}
\opt{D}{$6$}
\end{qoptions}
\qn{5} Let $A$ and $B$ be two non-singular real square matrices of order $4\times4$. If $AA^T=I=A^TA$ and $B=B^4$ such that $\det(A)\neq\det(B)$, then the value of $\displaystyle\sum_{r=1}^{20}\det\left(\mathrm{adj}\left(A^rB^{2r+1}\right)\right)$ is equal to -\\
(where $I$ is an identity matrix of order $4$ and $A^T$ represent transpose of $A$)
\begin{qoptions}
\opt{A}{$1$}
\opt{B}{$-1$}
\opt{C}{$0$}
\opt{D}{$4$}
\end{qoptions}
\qn{6} Let $\vec a,\vec b$ and $\vec c$ be $3$ mutually perpendicular unit vectors. If $\vec d$ is a unit vector which is equally inclined to $\vec a,\vec b$ and $\vec c$, then the value of $\left|\vec a+\vec b+\vec c+\vec d\right|^2$ is -
\begin{qoptions}
\opt{A}{$4+2\sqrt3$}
\opt{B}{$4+3\sqrt2$}
\opt{C}{$2+\sqrt5$}
\opt{D}{$\sqrt5-2$}
\end{qoptions}
\qn{7} Equation of lines\\
$L_1:\ 2x-2y+3z-2=0,\ x-y+z+1=0$\\
$L_2:\ x+2y-z-3=0,\ 3x-y+2z-1=0$\\
Distance of point $P(0,0,0)$ from the plane containing $L_1$ and $L_2$ and measured along the line $x=y=z$ is -
\begin{qoptions}
\opt{A}{$\dfrac1{3\sqrt2}$}
\opt{B}{$\dfrac{3\sqrt3}8$}
\opt{C}{$\dfrac{\sqrt3}8$}
\opt{D}{$\dfrac{\sqrt3}4$}
\end{qoptions}
\sect{SEC-2}
\qn{1} If $\alpha,\beta\in R$, $\beta\neq0$, $n\in N$ and $\displaystyle\lim_{n\to\infty}\dfrac{\sqrt[n^2]{(1!)(2!)(3!)\dots(n!)}}{n^\alpha}=\beta$, then
\begin{qoptions}
\opt{A}{$\alpha=\dfrac12$}
\opt{B}{$\beta=e^{-\frac34}$}
\opt{C}{$\alpha=1$}
\opt{D}{$\beta=\dfrac1{\sqrt e}$}
\end{qoptions}
\qn{2} If $y=f(x)$ is a solution of $\dfrac{dy}{d(x^2)}+\dfrac y{x^2}=\dfrac{\sin^{-1}x}{2x^2}$ and $f(0)=0$, then
\begin{qoptions}
\opt{A}{$f(1)=\dfrac\pi8$}
\opt{B}{$f(-1)=-\dfrac\pi8$}
\opt{C}{$f\!\left(\dfrac12\right)=\dfrac{\sqrt3}4-\dfrac\pi6$}
\opt{D}{$x^2f(x)=\dfrac12(2x^2-1)\sin^{-1}x+\dfrac14x\sqrt{1-x^2}+\dfrac14$}
\end{qoptions}
\qn{3} If $f(x)=\begin{vmatrix}5+\sin^2x&\cos^2x&4\sin2x\\ \sin^2x&5+\cos^2x&4\sin2x\\ \sin^2x&\cos^2x&5+4\sin2x\end{vmatrix}$, then which of the following is/are correct -
\begin{qoptions}
\opt{A}{$f(x)$ is a non-monotonic function}
\opt{B}{$f(x)=0$ has no real roots}
\opt{C}{$f'(x)$ has infinitely many maxima and minima}
\opt{D}{Number of integers in the range of $f(x)$ is $201$}
\end{qoptions}
\qn{4} If $\cos\dfrac{2\pi}{2019}\cdot\cos\dfrac{4\pi}{2019}\cdot\cos\dfrac{6\pi}{2019}\cdots\cos\dfrac{2018\pi}{2019}$ is $2^{-k}$ where $k=\dfrac pq$, $p,q\in N$ and $p$ and $q$ are relatively prime, then
\begin{qoptions}
\opt{A}{$(p+q)$ is a $4$-digit number}
\opt{B}{sum of the digits in $(p+q)$ is $2$}
\opt{C}{$(p+q)$ is a $3$-digit number}
\opt{D}{number of divisors of $p$ is $8$}
\end{qoptions}
\qn{5} For the curve $\sin x+\sin y=1$, $\displaystyle\lim_{x\to0}(x^\alpha)\dfrac{d^2y}{dx^2}$ exists then $\alpha$ can be
\begin{qoptions}
\opt{A}{$\dfrac32$}
\opt{B}{$2$}
\opt{C}{$\dfrac52$}
\opt{D}{$1$}
\end{qoptions}
\qn{6} A point $P$ moves inside a square with vertices $A(1,1)$, $B(-1,1)$, $C(-1,-1)$ and $D(1,-1)$ such that minimum $\{PA,PB,PC,PD\}\le1$. If the area bounded by the curve traced out by moving point $P$ is $\lambda$, then which of the following is/are correct -
\begin{qoptions}
\opt{A}{$\lambda$ is a rational number}
\opt{B}{$[\lambda]=3$ (where $[.]$ represents greatest integral function)}
\opt{C}{$\lambda$ is an irrational number}
\opt{D}{$[\lambda]=0$ (where $[.]$ represents greatest integral function)}
\end{qoptions}
\qn{7} $\displaystyle\lim_{x\to0}\dfrac{x^2\sin\beta x}{\tan\alpha x-\sin x}=4$ where $\alpha,\beta\in R$, then which of the following is/are correct -
\begin{qoptions}
\opt{A}{$\alpha=1$}
\opt{B}{$\beta=2$}
\opt{C}{$\alpha=2$}
\opt{D}{$\beta=1$}
\end{qoptions}
\sect{SEC-3}
\paraq{15}{16}
Let $\vec a,\vec b$ and $\vec c$ be three mutually perpendicular unit vectors
\qn{1} Let $\vec r$ be a unit vector satisfying
\[(\vec b-\vec c)\times(\vec r\times\vec a)+(\vec c-\vec a)\times(\vec r\times\vec b)+(\vec a-\vec b)\times(\vec r\times\vec c)=\vec 0\]
. Then $\vec r$ is equal to
\begin{qoptions}
\opt{A}{$\dfrac1{\sqrt3}(\vec a+\vec b+\vec c)$ only}
\opt{B}{$-\dfrac1{\sqrt3}(\vec a+\vec b+\vec c)$ only}
\opt{C}{Both option (A) and option (B) are correct}
\opt{D}{None of these}
\end{qoptions}
\qn{2} Let $\vec r$ be a vector having components $12,5$ and $13$ along $\vec a,\vec b$ and $\vec c$ respectively. If $\lambda$ is component of $\vec r$ along $(\vec a+\vec b+\vec c)$, then $[\lambda]=$\\
(where $[.]$ represents greatest integer function)
\begin{qoptions}
\opt{A}{$17$}
\opt{B}{$30$}
\opt{C}{$10$}
\opt{D}{$18$}
\end{qoptions}
\paraq{17}{18}
Let $f(x)$ and $g(x)$ be two quadratic polynomials $f(x)=ax^2+bx+c$, $g(x)=dx^2+ex+f$ where $a,b,c,d,e,f\in R$
\qn{3} If $a,b,c$ are the positive real numbers and $f(x)\ge0\ \forall\ x\in R$, then the minimum value of $\left(\dfrac{a+b+c}{b}\right)$ is
\begin{qoptions}
\opt{A}{$1$}
\opt{B}{$2$}
\opt{C}{$3$}
\opt{D}{$0$}
\end{qoptions}
\qn{4} If $H(x)=f(x).g(x)$ where $ad=1$ and $H(1)=1$, $H(2)=2$, $H(3)=3$, then $\left\{\dfrac{H(-1)+H(5)}{H(0)+H(4)}\right\}=$\\
(where $\{\ \}$ represents fractional function) is
\begin{qoptions}
\opt{A}{$\dfrac57$}
\opt{B}{$\dfrac27$}
\opt{C}{$\dfrac37$}
\opt{D}{$\dfrac47$}
\end{qoptions}
\end{multicols}
\end{document}
